\chapter{Conclusion and Future Work}\label{cap:conclfw}
This thesis set out to investigate whether a multi-user neurofeedback system could realistically be used in a classroom, not as a laboratory instrument but as a practical teaching aid. To answer that, a complete prototype was designed and built, covering the full path from EEG acquisition on per-student nodes, through real-time engagement estimation, to a teacher-facing application that brings live monitoring, lesson presentation, and post-lesson reporting together in one place.

The prototype was then put to work in a real classroom-style session with five participants, with the system recording everything on its own. It handled five simultaneous EEG streams for a full twenty-minute lesson without interruption, kept the dashboard responsive in near real time, and produced a clear per-slide and per-subject report afterwards. The engagement it reported behaved sensibly: it stayed high during the attentive opening of the lesson and reflected both the shared drop and the individual differences that appeared in the second half. On the central question of feasibility, the outcome is positive.

It is worth being clear about what this does and does not show. A single short session with five participants is far too small to draw conclusions about engagement or learning, and that was never the goal. What the work demonstrates is that an accessible, multi-user setup can turn raw EEG into organised, readable information that is tied to what was happening in the lesson, using hardware and software that a teacher could plausibly operate without laboratory conditions. The contribution of the thesis is therefore the working system itself, together with the evidence that the approach is practical, rather than any finding about how students learn.

Several directions follow naturally for the system. The engagement metric would benefit from being validated against an external reference instead of being trusted on its own, and the preprocessing could be strengthened to better handle the movement and contact artifacts that portable EEG is prone to. The interface could grow to support additional engagement measures, classroom grouping, and richer lesson analysis, while the deployment workflow could be made smoother so that larger numbers of acquisition nodes can be prepared quickly. Moving the prototype onto the cheaper single-electrode hardware it ultimately targets would also be a useful step, confirming in practice that the same results hold on the low-cost devices the project is meant for.

Beyond improving the system, the original motivation behind this work points to its most important continuation. From the very beginning, the intention was not only to build the tool but to use it to carry out a proper study of classroom engagement, on a scale where the results would actually mean something. With the prototype now shown to be feasible, that study becomes possible: a larger group of students, across several real lessons and over a longer period, would allow engagement patterns to be analysed seriously, compared between topics or teaching styles, and connected to the educational questions that a five-person session cannot answer. The system built here is intended precisely as the foundation for that kind of research.

In short, this thesis delivers a working, end-to-end neurofeedback system for classroom engagement monitoring and shows that it is feasible to run in practice. The natural next step is to take it from a feasibility prototype to a research instrument, and to use it for the larger, more rigorous study of classroom engagement that motivated the project in the first place.
